\subsection*{(2) Gravity --- the classical covariant field equation carries no free dimensionless Newton coupling under the named QFT/entanglement premises \mC; the $R^2$ inflation layer up to one truncation \mE/\mC}
The metric-sector \emph{field equation} itself carries no free dimensionless Newton coupling under those premises: fixed-\emph{volume} entanglement
equilibrium (Jacobson 2015) forces the \emph{full} covariant Einstein equation
$G_{ab}{+}\Lambda g_{ab}=\cthree^{-1}T_{ab}$ with \emph{both} coefficients fixed
($8\pi{=}1/\cthree$, $\Lambda$ from $\alpha$; \vref{v358}/\vref{v359}). What \emph{this} subsection details is
the complementary \emph{spectral-action / inflation} layer --- the $R+R^2$ readout --- whose closure is
\emph{not} a single hypothesis but a chain of three standard building blocks plus \emph{one} controlled
remainder:

\begin{center}
\begin{tikzpicture}[font=\footnotesize,
  st/.style={rectangle,rounded corners,draw=tfgreen!70!black,fill=tfgreen!8,thick,align=center,inner sep=3pt,minimum height=9mm},
  rest/.style={rectangle,rounded corners,draw=tfgold!80!black,fill=tfgold!12,very thick,align=center,inner sep=3pt,minimum height=9mm},
  ar/.style={-{Stealth[length=2mm]},thick,tfgray}]
  \node[st] (a) at (0,0) {spectral action\\(heat kernel $a_4$)\\$\to R{+}R^2{+}\operatorname{tr}\mathbb F^2$};
  \node[st] (b) at (3.7,0) {BRST/Hodge\\differential kills\\gauge directions};
  \node[st] (c) at (7.2,0) {FRW reduction\\$\to$ Starobinsky\\potential};
  \node[rest] (d) at (10.8,0) {\textbf{remainder:}\\low-curvature\\truncation\\$R^n\!\to\!R^2$};
  \draw[ar] (a)--(b); \draw[ar] (b)--(c); \draw[ar] (c)--(d);
\end{tikzpicture}
\end{center}

\begin{itemize}[leftmargin=1.4em]
\item \textbf{Spectral action:} the Seeley--DeWitt coefficients give Einstein--Hilbert
      ($a_2\!\sim\!-R/3$), $R^2$ ($a_4|_{R^2}\!\sim\!R^2/72$, coefficient $\tfrac1{6M^2}$) and gauge terms.
      Standard heat-kernel geometry (Gilkey), not a TFPT specificity; the scalaron ratio
      $M^2/\Mbar^2=6(4\pi)^2/f_0$ is fixed by the cutoff moment ($\cthree^7$). \mE\
      \veri{v36\_spectral\_action\_g2.py}
\item \textbf{Hodge/BRST closure:} the geometric BRST/Hodge differential quotients diffeomorphisms $+$
      local-frame gauge; the physical sector is its cohomology. A theorem as soon as the elliptic
      complex on the \emph{compactified} normal slice $S^2$ (APS boundary conditions) is
      Fredholm --- again standard index theory. \mE\ (given Fredholm)
\item \textbf{FRW reduction:} minisuperspace symmetry reduction onto the scale factor $+$ scalaron;
      yields the Starobinsky potential. Standard. \mE
\item \textbf{The only remainder:} the low-curvature \emph{truncation} --- higher curvature terms
      $R^n$ ($n\!\ge\!3$) reduce to the local $R^2$. A \emph{controlled approximation} with explicit
      domain of validity (curvature $\ll M^2$), \emph{not} a free assumption.
\end{itemize}

\par\smallskip\noindent\textbf{Sharpening (2): the truncation is explicitly bounded.}
$M_{\mathrm{Pl}}$ and $A_s=\tfrac{N_\star^2}{24\pi^2}c_3^7$ are determined \emph{exactly in the
low-curvature $R^2$ regime}. The truncation remainder is \emph{not} uncontrolled: on the Starobinsky
attractor during observable inflation $R/M^2\sim\tfrac1{2N_\star}$, i.e.\ an expansion parameter
$\approx0.009$ at $N_\star\!\approx\!55$. The $R^{n\ge3}$ corrections to $(n_s,A_s)$ are therefore
$O(1/N_\star^2)\approx3\times10^{-4}$ ($\sim0.03\%$) --- explicitly bounded and tiny. The leading
predictions $n_s=1-2/N_\star$ and $r=12/N_\star^2$ (the band over $N_\star\in[50,60]$) hit Planck. The
status is thus not ``conditional on an unproved Hodge theorem'' but ``exact up to an
$O(1/N_\star^2)$ truncation''. \textbf{Remaining path:} the Fredholm property of the seam elliptic
complex (APS, standard index theory). \mE\ in the regime.

\subsection*{(3) Strong CP --- decoupled from the mass-gap problem \mE/\mC}
\textbf{The key reduction:} $\theta_{\mathrm{eff}}=0$ needs \emph{neither} the full OS package \emph{nor}
the mass gap. It follows from three \emph{purely structural} facts plus reflection positivity on the
admissible branch:

\begin{center}
\begin{tikzpicture}[font=\footnotesize,
  box/.style={rectangle,rounded corners,thick,align=center,inner sep=3pt,minimum height=8mm},
  ar/.style={-{Stealth[length=2mm]},thick,tfgray}]
  \node[box,draw=tfgreen!70!black,fill=tfgreen!8] (i)   at (0,1.0) {(i) polar structure\\$M_fP_R{+}M_f^\dagger P_L$};
  \node[box,draw=tfgreen!70!black,fill=tfgreen!8] (ii)  at (0,0.0) {(ii) $\gamma_5$-Hermiticity\\$\Rightarrow \det$ real};
  \node[box,draw=tfgreen!70!black,fill=tfgreen!8] (iii) at (0,-1.0){(iii) sheet involution\\(antiunitary, CP)};
  \node[box,draw=tfblue,fill=tfblue!8] (rp) at (4.0,0) {$+$ reflection\\positivity (theorem)};
  \node[box,draw=tfgold!80!black,fill=tfgold!12,very thick] (cp) at (7.6,0) {$\boxed{\theta_{\mathrm{eff}}=0}$\\\mE\ (given RP)};
  \draw[ar] (i)--(rp); \draw[ar] (ii)--(rp); \draw[ar] (iii)--(rp); \draw[ar] (rp)--(cp);
\end{tikzpicture}
\end{center}

\noindent The three facts are algebraic and proved on $\operatorname{Ran}(P_{\mathrm{adm}})$ (the
admissibility three-sector motif). \textbf{Separately}, the admissibility / OS-reconstruction layer builds the \emph{full} OS
closure as its own chain: admissible reflection positivity (theorem) $+$ mass-gap stability in the
thermodynamic limit (Block 4, ``parent domination'') $+$ clustering $\Rightarrow$ OS reconstruction
(theorem). Reflection positivity, covariance, clustering and OS reconstruction are there \emph{theorems
given the gap}; the only genuinely hard input is the \textbf{mass-gap stability} (parent-domination
lemma).

\par\smallskip\noindent\textbf{Sharpening (3).}
Strong CP is \textbf{decoupled} from the mass-gap problem: $\theta_{\mathrm{eff}}=0$ is \mE\ given
reflection positivity (three structural facts). Only the \emph{full} QFT dynamics rests on mass-gap
stability --- the remaining constructive \mC. \textbf{Hardening path:} prove the parent-domination
lemma rigorously; all other OS axioms then follow.

\par\smallskip\noindent\textbf{CAR/Pfaffian form (\veri{v173\_pfaffian\_cp\_car.py}, \mE).} The same chain
is a Pfaffian statement in the self-dual $16$-Majorana CAR model: the sheet-odd Calder\'on involution
$\{D_\Sigma,\Sigma\}=0$ pairs the seam Dirac spectrum into $(\pm\lambda)$; $\gamma_5$-Hermiticity
($D^\dagger=\gamma_5 D\gamma_5$) makes the fermion-measure Pfaffian real, so $\arg\mathrm{Pf}\in\{0,\pi\}$;
and a $\theta$-term gives $Z(\theta)=e^{i\theta}|\mathrm{Pf}|$, whose reflection positivity ($Z>0$)
excludes $\theta=\pi$ ($\operatorname{Re}Z<0$) and every $\theta\notin\{0,\pi\}$ ($\operatorname{Im}Z\neq0$).
So $\theta_{\mathrm{eff}}=0$ is a topological null theorem --- the antisymmetric Calder\'on pairing makes
the Pfaffian real and RP picks the positive branch (machine-checked model demonstration).

\subsection*{(4) Constructive QFT closure of the admissible physical sector \mE/\mC}
This is the constructive QFT closure of the \emph{admissible physical sector} --- \emph{not} the
unprojected ambient theory: the unprojected ambient metric measure is not constructed here and remains the
$G_{\mathrm{metric}}$ frontier item (\S ambient QG). On the admissible sector, the input previously listed
as ``open'' --- the parent-domination / mass-gap lemma from (3) --- is \emph{not} analytically open: it is
\emph{computed} by the \textbf{explicit discrete spectrum} of the transport operator.

\begin{keybox}[The mass-gap input is explicit (verified)]
The admissible transport is the finite cyclic operator $D_y=y\mathbf 1-\delta_{\mathrm{ph}}U_6$ on $C_6$.
Its correlation/transfer operator has the \emph{closed} spectrum (the cusp cubic
$P(z)=(z-1)(z-\tfrac{64}{729})(z-\tfrac1{729})$):
\[
  \mathrm{spec}=\Bigl\{\bigl(\tfrac{k}{3}\bigr)^6:k=3,2,1\Bigr\}=\Bigl\{1,\ \bigl(\tfrac23\bigr)^6,\ \bigl(\tfrac13\bigr)^6\Bigr\},
  \qquad
  \boxed{\ \Delta_{\mathrm{gap}}=-\log\bigl(\tfrac23\bigr)^6=6\log\tfrac32=2.4328\ }
\]
The mass gap is thus \emph{positive, explicit and a compiler number} ($\tfrac32$ is the carrier rank
ratio $B$). \mE
\end{keybox}

\begin{center}
\begin{tikzpicture}[font=\footnotesize]
  \draw[->,thick,tfgray] (0,-0.2)--(0,3.3) node[above]{$-\log\lambda$};
  \foreach \y/\lab in {0/{$\lambda_1{=}1$ (vacuum)}, 1.46/{$\lambda_2{=}(2/3)^6$}, 3.0/{$\lambda_3{=}(1/3)^6$}}{
    \draw[thick,tfblue] (-0.15,\y)--(1.6,\y); \node[anchor=west,tfblue] at (1.7,\y) {\lab};}
  \draw[{Stealth[length=2mm]}-{Stealth[length=2mm]},very thick,tfgold!80!black] (0.7,0.05)--(0.7,1.41);
  \node[anchor=west,tfgold!80!black] at (0.75,0.73) {$\Delta_{\mathrm{gap}}=6\log\tfrac32$};
\end{tikzpicture}
\end{center}

\noindent\textbf{The closure chain} (now complete):
\begin{enumerate}[leftmargin=1.6em]
\item \emph{Explicit gap.} The finite cyclic transfer operator is gapped with $\Delta=6\log\tfrac32$ ---
      this is the ``parent gapped quasilocal'' object of parent domination, written out. \mE
\item \emph{Thermodynamic limit.} A \emph{uniform} gap of a finite-range transfer operator gives, by the
      standard cluster expansion, exponential clustering and stability in the limit. \mE\ (standard)
\item \emph{OS package.} With admissible reflection positivity (the OS-reconstruction theorem) $+$ gap $+$ clustering
      all Osterwalder--Schrader axioms hold $\Rightarrow$ OS reconstruction: Hilbert space, Hamiltonian
      and a \emph{measure} on the admissible sector. \mE
\item \emph{Bridge values $=$ correlators.} The fixed values (masses, mixings, $\alpha$) are the
      Schwinger functions of this reconstructed measure.
\item \emph{Gauge directions trivial.} The unphysical/off-shell directions are cohomologically dead via
      the BRST/Hodge differential (question 2) --- so on the \emph{admissible gapped IR sector} this is the
      full physical dynamics; the \emph{ambient} metric/projective measure is the separate $G_{\mathrm{metric}}$
      gate (below), not claimed here.
\end{enumerate}

\par\smallskip\noindent\textbf{Sharpening (4): what is and is not solved} (\mE/\mC).
The \emph{admissible gapped transfer sector} is OS-reconstructed and its measure exists because the gap is
\emph{explicit} --- \emph{under the stated hypotheses}: locality, reflection positivity after integration, a
\emph{uniform} gap in the regulator limit, tightness, and the Ward/BRST gauge quotient. Within those, point
(4) is \mE\ on the admissible sector and the only standard step is gap-$\to$-limit stability (cluster
expansion). The \emph{perturbative} 4D readout built on this sector now closes \emph{to all orders} --- read strictly as a
\emph{local perturbative construction under the stated Epstein--Glaser/BRST hypotheses} (the massless-sector weak
adiabatic limit is the cited Duch theorem, \texttt{arXiv:1801.10147}, whose hypotheses remain to be verified for the
concrete TFPT model; no confinement, mass gap, asymptotic charged states, nonperturbative Higgs vacuum or controlled
full S-matrix follows, \texttt{QFT4D.OS.RECON.01} in \texttt{tfpt\_research\_contracts}) --- as a typed
Epstein--Glaser/BRST contract (\veri{v381\_qft4d\_eg\_allorder.py}, \texttt{QFT4D.EG.ALLORDER.01}): dimension-4
power-counting gives a finite counterterm space, BRST nilpotency $s^2{=}0$ holds for the carrier $su(3){\times}su(2)$,
and the seam gap gives the adiabatic limit, with all-order $T_n$ existence (causal factorisation) and the
Slavnov--Taylor identity imported --- matter$+$gauge only (the $R^2$/Weyl$^2$ ghost is the resummation
\vref{v304}/\vref{v370}/\vref{v380}, not a local vertex). \textbf{What is \emph{not} constructed here:} the full $4$D continuum, whose cited existence theorems
stay \mO\ (the seam being \emph{closed modulo cited theorems} on its MMST route \texttt{SEAM.EQUIV.MMST.01}, \vref{v336}; the parent \texttt{SEAM.EQUIV.01} stays \mO), and the ambient metric/projective QG
measure, now a \mC\ redundancy (\vref{v369}; G6 projective limit, \texttt{tfpt\_research\_contracts}).
So: \emph{admissible gapped transfer sector solved under RP/gap hypotheses; the seam closed modulo cited theorems on
its MMST route (\texttt{SEAM.EQUIV.MMST.01}; parent \mO);
ambient metric/projective QG measure a \mC\ redundancy}.
\par\smallskip\noindent\textbf{Three terms, kept apart} (to read ``QFT closure'' correctly): $\mathrm{QFT}_{\mathrm{adm}}$
is the gapped admissible transfer theory with OS reconstruction (closed under RP and gap, \mE);
$\mathrm{AQFT}_{\mathrm{bdry}}$ is the holomorphic $c=8$ $(E_8)_1$ boundary net (now \emph{closed modulo cited
theorems} on the MMST route \texttt{SEAM.EQUIV.MMST.01} via the lattice/$S_3$ stack and \texttt{FORM.SEAM.MMST.01}, \mC); and $\mathrm{QG}_{\mathrm{amb}}$ is the
unprojected ambient metric projective measure (now a \mC\ redundancy, \vref{v369}). TFPT constructs $\mathrm{QFT}_{\mathrm{adm}}$ and
$\mathrm{AQFT}_{\mathrm{bdry}}$; it does \emph{not} construct $\mathrm{QG}_{\mathrm{amb}}$. Whenever ``QFT
closure'' appears, it means the first two, never the third.
\par\smallskip\noindent\textbf{Since assembled into one relative object (Modular Spectral Closure).} The two
constructed terms $\mathrm{QFT}_{\mathrm{adm}}+\mathrm{AQFT}_{\mathrm{bdry}}$ are now a \emph{single} relative
datum $\mathsf{TFPT}_{\mathrm{QFT}}=(\mathcal A_\Sigma,\omega_\Sigma,\Delta_\Sigma,\rho,A_F,H_F,D_F,J,\gamma,
S_{\mathrm{rel}})$: the explicit $96$-dim finite spectral triple is built (\veri{v252\_full\_finite\_triple.py}),
its Dirac is the covariance induction of the seam KMS state (\veri{v258\_dirac\_covariance\_induction.py}), the
spectral-action cutoff is that KMS weight ($f_2/f_0{=}1$, \veri{v259\_modular\_cutoff\_kappa.py}), and seam,
carrier-$16$ and $E_8$ sit on one Kummer/K3 (\veri{v260\_k3\_kummer\_unification.py}), the whole assembly
certified together (\veri{v261\_modular\_spectral\_closure.py}) --- closed modulo the single premise
\texttt{QGEO.SYM.01}. \emph{What this still does \emph{not} close} is the dimensionful/ambient remainder: the
full measured pole masses ($\times\,v_{\mathrm{geo}}$ $+$ standard RG, \mC), the absolute QCD scales
($m_p/m_e$, now a typed \mC\ transfer, \vref{v374}), the complete PMNS dynamics (\mC/\mO), and the real $4$D
interacting QFT $\mathrm{QG}_{\mathrm{amb}}$ (a \mC\ redundancy, \vref{v369}; only the cited continuum existence
stays \mO, \vref{v336}).

\subsection*{(5) Axioms P1/P2 --- P2 machine-verified, P1 as an interface \mE/\mO}
Here the hardening is \emph{empirical} and concrete. The Lean 4 proof (\texttt{lean4-carrier-rigidity},
toolchain \texttt{v4.29.1} $+$ Mathlib) builds cleanly:
\begin{center}\small
\renewcommand{\arraystretch}{1.15}
\begin{tabularx}{\linewidth}{@{}lY@{}}
\toprule
\textbf{Verified (Lean, 0 \texttt{sorry})} & \textbf{Content} \\
\midrule
\texttt{Polarization.sixY\_carrier\_polynomial} & $6Y^2-Y-1=0$ as a corollary of orthogonal idempotents \\
\texttt{Rigidity.unique\_carrier\_pair} & $(q_-,q_+)=(-2,3)$ unique (integer rigidity) \\
\texttt{Hypercharge.trace\_Y}, \texttt{\_carrier\_polynomial} & $Y=\operatorname{diag}(-\tfrac13,-\tfrac13,-\tfrac13,\tfrac12,\tfrac12)$, $\operatorname{tr}Y=0$ \\
\texttt{CarrierData.CarrierPremises.*} & the same statements \emph{abstractly} over any field/module \\
\texttt{YukawaStageDExistence}, \texttt{YukawaTrilinearForm} & $\dim E=3$ (family number) from Yukawa stage D \\
\bottomrule
\end{tabularx}
\end{center}

\par\smallskip\noindent\textbf{Sharpening (5): build result this round (reproducible).}
\emph{Provenance:} canonical shipped folder \texttt{lean4-carrier-rigidity/} (source repo: \texttt{experiments/lean4-carrier-rigidity/}); toolchain
\texttt{leanprover/lean4:\vref{v4}.29.1}; one command runs the hard gate, \texttt{bash scripts/audit.sh},
which on this round returns (re-run, $\sim7$\,s with cached \texttt{.lake/build}):
\texttt{lake build} $\Rightarrow$ success (1886 jobs); \textbf{0} \texttt{sorry}/\texttt{admit}; no
\texttt{axiom}/\texttt{constant}/\texttt{unsafe}/\texttt{opaque}/\texttt{partial} declarations; and
\texttt{\#print axioms} on every headline theorem reports \emph{only} the three kernel axioms
(\texttt{propext}, \texttt{Classical.choice}, \texttt{Quot.sound}) --- no domain-specific axioms, no
\texttt{sorry}-laundering. So the formalised marker is reproducible, not a textual claim.

\noindent\textbf{P1/P2 balance:} the \emph{algebra} of P2 (carrier $\to$ hypercharge, rigidity, family
number) is thereby \emph{freed from axiom status} --- it is a machine-checked theorem from the
idempotent axioms. P1 (boundary kernel, Calder\'on/seam-winding) is formalised as \emph{typed
interfaces} (\texttt{BoundaryYukawaKernel}, \texttt{CalderonInterface}, \texttt{SeamWinding}): its
analytic premises are imported hypotheses, not derived. \textbf{Irreducible remainder:} ``two orthogonal
idempotents $+$ the analytic Calder\'on/winding interface''. \textbf{Hardening path:} formalise the
index-theoretic premises directly (substantially more work, beyond the current Lean scope). \mO

\begin{gapbox}[Overall balance of the five questions]
\renewcommand{\arraystretch}{1.2}\small
\begin{tabularx}{\linewidth}{@{}p{0.5cm}p{4.3cm}Y@{}}
\toprule
\# & \textbf{Status now} & \textbf{Isolated remainder} \\
\midrule
1 & \mE\ (closed) & unitarity of $\nabla_F^\star$ (parabolic Hodge structure) \\
2 & \mE\ in the $R^2$ regime & controlled $R^{n\ge3}$ truncation $+$ Fredholm \\
3 & \mE\ (strong CP, given RP) & full dynamics: admissible gapped sector via (4) under RP/gap; ambient QG a \mC\ redundancy (\vref{v369}) \\
4 & \mE\ adm.\ sector & admissible-sector measure closed (gap $6\log\tfrac32$, OS) \emph{under RP/gap hyp.}; ambient metric measure a \mC\ redundancy ($G_{\mathrm{metric}}$, \vref{v369}) \\
5 & \mE\ (P2 algebra, Lean) & P1 analytics as an interface; index premises \\
\bottomrule
\end{tabularx}
\markerkey

\smallskip
\noindent The five points are \emph{closed, reduced or typed}: (1) reduced to the
$(U_{\mathrm{wall}})$ selector (Mehta--Seshadri framework $+$ stable-point selection), (3) decoupled from the
mass-gap problem, (4) the admissible transfer model closed by the explicit discrete gap
$6\log\tfrac32$ (continuum/metric conditional), (5) P2 machine-verified; (2) the local field equation is
\emph{parameter-free} (full covariant $G_{ab}{+}\Lambda g_{ab}{=}\cthree^{-1}T_{ab}$, \vref{v358}/\vref{v359}),
its $R^2$ inflation layer exact up to a controlled truncation. What remains are \emph{standard steps} (cluster expansion, Fredholm, truncation
bound), the \emph{two declared inputs} P1/P2, and the two explicitly typed analytic gates
$(U_{\mathrm{wall}})$ and $G_{\mathrm{metric}}$ --- no \emph{hidden} finite arithmetic gap in the compiled
discrete core.
\end{gapbox}

% =====================================================================
\section{How the original-paper predictions simplify}
% =====================================================================
The same numbers, now as \emph{closed} expressions instead of distributed
theorems:

\begin{center}\small
\renewcommand{\arraystretch}{1.25}
\begin{tabularx}{\linewidth}{@{}p{3.2cm}p{4.4cm}Y@{}}
\toprule
\textbf{Prediction} & \textbf{Original paper (distributed)} & \textbf{Compiler form (one line)} \\
\midrule
$\alpha^{-1}$ & transport budget $\sum L{+}N_\Phi{=}41$, multi-stage & unique root of $F_{U(1)}(\alpha){=}0$, $137.0359992$ \\
fermion masses & per-fermion Yukawa $+$ $\Lambda$ table & \emph{one} $\phiz$-ladder: $\hat m{=}\tfrac{\vgeo\pi}{\sqrt2}c\,(\phiz)^{k}$ \\
lepton ratios & separate resolvent evaluations & $\tfrac87\phiz,\ \tfrac{12}7(\phiz)^2,\ \tfrac32\phiz$ \\
CKM & holonomy matrix, theorems & $s_{12}{=}\lambda_C,\ s_{23}{=}\tfrac{\phiz}{1+\lambda_C},\ s_{13}{=}\tfrac{\lambda_C^3}3$ \\
$\sin^2\theta_{13}$ & neutrino-closure readout & $\phiz\,e^{-5/6}$ \\
$\sin^2\theta_{12}$ & not stated & seed (freeze) $\tfrac13-\tfrac{\phiz}{2}{=}0.306747$ ($\varepsilon{=}q(A_3)\phiz$) \\
flavor matrix & transport theorem (audit) & residue matrix, $\det{=}h(D_5)$, minors $2{\cdot}3{\cdot}5{=}h(E_8)$ \\
$\Lambda$ (cosmol.) & determinant line $e^{-2\alpha^{-1}}$ & $\Lambda\sim e^{-2\alpha^{-1}}$ (engine) \\
$A_s,n_s,r$ & $R^2$ branch $+$ reheating & $A_s{=}\tfrac{N_\star^2}{24\pi^2}c_3^7,\ n_s{=}1{-}\tfrac2{N_\star},\ r{=}\tfrac{12}{N_\star^2}$ \\
strong CP / gap & Block-4 mass gap (assumption) & explicit gap $6\log\tfrac32$ (cusp spectrum) \\
\bottomrule
\end{tabularx}
\end{center}

\par\smallskip\noindent\textbf{What improves about the predictions.}
\emph{(i) From theorems to formulas:} quantities that arose through multi-stage theorems
($\alpha^{-1}$, flavor transport, masses) are now \emph{one-line} closed expressions.
\emph{(ii) Fewer inputs:} the whole mass spectrum hangs on \emph{one} seed $\phiz$ and \emph{one}
matrix --- the $\Lambda$ table (formerly 9 numbers) becomes the $\phiz$-ladder.
\emph{(iii) New sharp tests:} residue invariants ($\det 8$, minors $2{,}3{,}5$), the explicit gap
$6\log\tfrac32$, and the inflation band ($n_s,r$ over $N_\star\in[50,60]$) --- all fit-free and
independently checkable.

% =====================================================================
\section*{Conclusion}
% =====================================================================
\addcontentsline{toc}{section}{Conclusion}
The simplification is substantial: the discrete Standard-Model \emph{structure} (charges, families, flavor operators) is derived under the axioms; $4$D existence, chiral measure, mirror gap, couplings and neutrino dynamics are open.  The $\phiz$-ladder formula $\hat m=\tfrac{\vgeo}{\sqrt2}\lamY^{L}\Lambda$ with
\emph{one} seed $\phiz$ and \emph{one} fixed residue matrix whose spectral invariants
($\Nfam^2$, $2{\cdot}3{\cdot}5{=}h(E_8)$, $h(D_5)$) are pure compiler numbers. Through the $\phiz$-ladder all
nine word-lengths are fixed; the \emph{charged-lepton} amplitudes are closed, and the \emph{quark} mass
\emph{ratios} are closed too (integer Plücker readouts on the derived selector stratum, \vref{v49}/\vref{v71}) --- only the
\emph{absolute} quark amplitude scale reduces to the $(U_{\mathrm{wall}})$ anchor. On
the SM side the remainder is the \emph{conditional} status of $\theta_{12}$ (reduced to the
seam-misalignment lemma), the absolute quark amplitude scale, the full PMNS completion, and the scheme-RG
dimensionful masses ($m_H$ rests on the closed UV quartic).

On the five open questions the progress is equally concrete: (1) the flavor residual is \emph{reduced} to
the $(U_{\mathrm{wall}})$ selector (Mehta--Seshadri framework $+$ stable-point selection); (3) strong CP is
\emph{decoupled} from the mass-gap problem; (4) the \emph{admissible
bridge-sector} measure is closed (reflection positivity from the cusp spectrum, OS reconstruction, gap
$6\log\tfrac32$) --- \emph{not} the full ambient $4$D dynamics, whose unprojected measure remains
conditional on ambient RP; (5) the P2 algebra is \emph{machine-verified} (Lean, 1886 jobs, 0
\texttt{sorry}, clean axioms); (2) gravity's local field equation is \emph{parameter-free} (full covariant
$G_{ab}{+}\Lambda g_{ab}{=}\cthree^{-1}T_{ab}$, \vref{v358}/\vref{v359}), the $R^2$ inflation layer exact up to
a controlled truncation. What remains are the cited continuum-existence step (the seam closed modulo cited
theorems on its MMST route \texttt{SEAM.EQUIV.MMST.01}; the ambient measure a \mC\ redundancy, \vref{v369}) and the two declared inputs P1/P2 --- a
sharply outlined, short remainder.

\clearpage
\part{The flavor block: parabolic weights and the transport theorem}
\begin{abstract}
This note closes the last structural gap of the $D_5\times A_3$ compiler by deriving the
charged-fermion transport word-lengths $L_u=(7,3,0)$, $L_d=(7,5,2)$, $L_e=(8,5,3)$ --- equivalently the
residue triple $\{0,1,3\}$ and its $D_6$ orbit --- from first principles, rather than from the legacy
resolvent construction. The correct framework is parabolic Higgs/bundle theory on the rigid family curve
$X_f^\circ\cong\mathbb P^1\setminus\mu_4$, with the three transport cusps $\{1,\tfrac23,\tfrac13\}$
promoted to parabolic weights $\{0,\tfrac13,\tfrac23\}$ at each of the four punctures. We prove the
exact scaffolding (parabolic degree zero forces $\deg E=-|\mu_4|=-4$; the word-length splits as
$L=6n+r$ with winding $n=$ line-bundle degree and residue $r\in\mathbb Z_6$), record what the
Riemann--Hilbert side fixes (a rigid $SU(3)$ system whose discrete rotation is the $A_3$ Coxeter
element), and --- the main new result --- \emph{derive} the anchor multiset: the $\mathbb Z_4$
symmetrisation of the Fuchsian residues collapses to $\sum_k R^kA_0R^{-k}=4\operatorname{diag}(A_0)$, so
the exponents at infinity (the splitting type) are integers summing to $4$; Schur--Horn plus parabolic
stability force them to be $\{1,1,2\}$, exactly the per-sector anchors $(a_u,a_d,a_e)$ and the degrees of
$\mathcal O(-2)\oplus\mathcal O(-1)^2$. The per-sector \emph{assignment} then follows from the $D_5$ cusp
ordering via stability (the largest cusp, lepton, carries the most-negative summand $\mathcal O(-2)$,
i.e.\ the exponent $2$): $(a_u,a_d,a_e)=(1,1,2)$ is fully fixed. The anchor vector itself is a tiny
compiler --- $\|a\|_1=4=|\mu_4|=-\deg E$, $\|a\|_2^2=6=|R^+(A_3)|$, $\prod a_i=2$ (sheet) --- and the
whole integer landscape $40,41,48,60,120,240,248,128$ collapses to $d_p:=-\deg E=4$ times a Coxeter
datum, with the scalaron exponent $7=-\deg E+\operatorname{rk}E=4+3$. Finally the six non-anchor residues
are shown to be \emph{not free} either: the residue sets are the $D_6$ orbit of the unique
distinct-distance triple $\{0,1,3\}$ (pinned by their sums $(4,8,10)=(|\mu_4|,h(D_5),\AL)$), and removing
the anchor and ordering by hierarchy reproduces the full matrix $\{(7,3,0),(7,5,2),(8,5,3)\}$ exactly.
This is packaged as the \textbf{flavor transport theorem} (Thm.~\ref{thm:flavor}): under three
hypotheses (H1) non-degeneracy, (H2) cusp $D_6$-labelling, (H3) anchor$+$winding --- of which (H3) is
derived and (H1) is a genericity statement --- the full matrix $\{(7,3,0),(7,5,2),(8,5,3)\}$ is forced,
with the combinatorial core proved by complete enumeration. (H2)'s answer is moreover \emph{not} free:
The transport-kernel theorem derives the down branch $\{1,2,5\}$ (the spectral-selector branch,
$\det R=h(D_5)$, principal minors $(2,3,5)$) from the uniquely-fixed family holonomy; only the
parabolic$\leftrightarrow$transport equivalence remains open. Status flags: \mE\ exact identity, \mE\
lattice/Lie theorem, \mE\ numerical, \mC\ conditional/project, \mO\ open.
\end{abstract}

\begin{goal}
Derive the integer matrix
\[
  L=\begin{pmatrix}7&3&0\\7&5&2\\8&5&3\end{pmatrix}
  \quad(\text{rows }u,d,e;\ \text{columns gen }1,2,3),\qquad \sum_{f,j}L_{f,j}=40 ,
\]
as parabolic degrees of a rigid parabolic bundle on $\mathbb P^1\setminus\mu_4$, closing the last
flavor input of the compiler.
\end{goal}

% =====================================================================
